\documentclass[12pt, a4paper]{article}
\usepackage[utf8]{inputenc}
\usepackage[T1]{fontenc}
\usepackage[french]{babel}
\usepackage{amsmath, amssymb}
\usepackage{geometry}
\usepackage{array}
\usepackage{listings}
\geometry{margin=2.5cm}

\lstset{basicstyle=\ttfamily\small, breaklines=true, frame=single,
  columns=fullflexible, keepspaces=true}

\title{MN2 -- TP Scilab\\Proposition de correction}
% Pour plusieurs auteurs : \author{Nom 1 \and Nom 2 \and Nom 3}
\author{Zonou Zinta-Kevin}
\date{\today}

\begin{document}
\maketitle

\noindent\textbf{Vérification :} toutes les valeurs numériques ci-dessous ont été recalculées avec Scilab 2024.0.0 (routines de l'Annexe A et blocs TODO de l'Annexe B), et confirment les résultats obtenus précédemment en Python.

%----------------------------------------------------------------
\section{Contrôle sur $f(x)=e^x$}

Référence : $I_{\mathrm{ref}}=\displaystyle\int_0^1 e^x\,dx=e-1\approx 1{,}718281828$.

\begin{center}
\small
\begin{tabular}{|r|c|c|c|c|c|c|}
\hline
$n$ & $T_n$ & $S_n$ & $E_n(T)$ & $E_n(S)$ & $p(T)$ & $p(S)$\\
\hline
4  & 1,727221905 & 1,718318842 & $8{,}940\cdot10^{-3}$ & $3{,}701\cdot10^{-5}$ & 1,999 & 3,992\\
8  & 1,720518592 & 1,718284155 & $2{,}237\cdot10^{-3}$ & $2{,}326\cdot10^{-6}$ & 2,000 & 3,998\\
16 & 1,718841129 & 1,718281974 & $5{,}593\cdot10^{-4}$ & $1{,}456\cdot10^{-7}$ & 2,000 & 3,999\\
32 & 1,718421660 & 1,718281838 & $1{,}398\cdot10^{-4}$ & $9{,}103\cdot10^{-9}$ & -- & --\\
\hline
\end{tabular}
\end{center}

\noindent\textbf{Question de décision.} Oui : les ordres observés tendent vers $2$ (trapèzes) et $4$ (Simpson), comme prévu par le CM1.
Le calcul de $p=\log(E_n/E_{2n})/\log 2$ n'est fiable que si les erreurs sont nettement supérieures à l'erreur d'arrondi (de l'ordre de $10^{-16}$). Si l'erreur est nulle, le rapport est indéfini ; si elle est proche de l'arrondi, elle est dominée par le bruit numérique. D'où le retour \texttt{\%nan} de \texttt{ordre\_observe} lorsque $E\le 0$.

%----------------------------------------------------------------
\section{Intégrale discrète : la pompe}

Temps en heures : $0,\ \tfrac16,\ \tfrac13,\ \tfrac12,\ \tfrac23$ (pas $h=\tfrac16$ h), $P=2{,}0;\ 2{,}6;\ 3{,}0;\ 2{,}5;\ 1{,}9$ kW.

\[
T_4 = 1{,}675000\ \text{kWh},\qquad
S_4 = 1{,}683333\ \text{kWh},\qquad
|T_4-S_4| = 0{,}008333\ \text{kWh}.
\]

Borne due aux lectures : $0{,}1\times 40/60 = 0{,}066667$ kWh. L'écart entre méthodes est environ huit fois plus petit que cette contribution maximale.

\medskip
\noindent\textbf{Question d'ingénieur.} On communiquerait une énergie d'environ $1{,}68$ kWh sur 40 min, avec une incertitude d'environ $\pm 0{,}07$ kWh ($\approx 4\,\%$). Le choix trapèzes/Simpson est négligeable devant l'incertitude des mesures. Avec cinq relevés espacés de dix minutes, il n'existe ni maillages emboîtés ni référence exacte : on ne peut donc pas vérifier la convergence ni calculer un ordre.

%----------------------------------------------------------------
\section{Contrôles et remise}

Pour $P(t)=10+4t-t^2$ aux temps $0,1,2,3,4$ s : $T_4=50$ J et $S_4=152/3\approx 50{,}667$ J. Les routines sont validées. Si les temps sont irréguliers, seuls les trapèzes sont utilisables ; Simpson uniforme doit refuser ces données.

%----------------------------------------------------------------
\section*{Annexe B -- Blocs TODO complétés}

\begin{lstlisting}
// TODO 1
ET(k) = trapezes(t, y);
ES(k) = simpson(t, y);
// TODO 2
errT(k) = abs(ET(k) - Iref);
errS(k) = abs(ES(k) - Iref);
// TODO 3
pT = ordre_observe(errT(k), errT(k+1));
pS = ordre_observe(errS(k), errS(k+1));
// TODO 4
th(i) = tmin(i)/60;
// TODO 5
Etrap = trapezes(th, Pkw);
Esimp = simpson(th, Pkw);
\end{lstlisting}

\end{document}
